% ============================================================
% §5  Dark Energy from Fluctuations of Generated Space
% ============================================================
\section{Dark energy}\label{sec:ciemna}

In standard cosmology dark energy is a cosmological constant~$\Lambda$
inserted by hand into the Einstein equations, or a quintessence field
with an arbitrary potential.  In TGP dark energy arises from the
\textbf{dynamics of generated space itself}: condensation of matter
amplifies fluctuations of the field~$\Phi$, which generate additional
space.

% ----------------------------------------------------------
\subsection{Mechanism}\label{ssec:de-mech}

\begin{proposition}[Dark energy from $\Phi$~fluctuations]%
\label{prop:DE}
The field~$\Phi$ undergoes fluctuations---both quantum (on small
scales) and classical/turbulent (superposition of contributions from
many sources on large scales).  These fluctuations generate
\textbf{additional effective space}, because the large-scale average
$\langle\Phi\rangle$ is elevated relative to the value that the same
mass would produce if distributed homogeneously.
\end{proposition}

The detailed mechanism proceeds in four steps:
\begin{enumerate}[nosep]
  \item Every particle generates a field~$\Phi$ in its vicinity.
  \item When particles condense into structures (galaxies, clusters),
        their $\Phi$~fields superpose \textbf{nonlinearly}.
  \item The nonlinearity $\mathcal{N}[\Phi]$ implies that the
        average $\langle\Phi\rangle$ in a structured region is
        \textbf{higher} than~$\Phi$ for the same mass spread
        homogeneously:
        \begin{equation}\label{eq:Phi-excess}
            \langle\Phi_{\mathrm{struct}}\rangle
            > \Phi_{\mathrm{homogeneous}}.
        \end{equation}
  \item The excess
        $\Delta\Phi = \langle\Phi_{\mathrm{struct}}\rangle
        - \Phi_{\mathrm{homogeneous}}$
        represents \textbf{additional space}---space that ``should
        not exist'' if matter were distributed uniformly.
        This additional space manifests as \textbf{accelerated
        expansion}: there is more ``room'' than matter alone
        would warrant, so the Universe expands faster than a
        matter-only model predicts.
\end{enumerate}

\begin{remark}[Flattening the problem]
In standard cosmology one asks: ``Why does $\Lambda$ have
\emph{this} particular value?'' (fine-tuning problem).
In TGP the question becomes: ``How much additional space do
$\Phi$~fluctuations generate for a given matter distribution?''
---the answer follows from dynamics, not from parameter tuning.
\end{remark}

% ----------------------------------------------------------
\subsection{Link to structure growth}\label{ssec:de-struktury}

A key prediction: dark energy in TGP \textbf{grows} as cosmic
structures form.
\begin{itemize}[nosep]
  \item In the early Universe matter is nearly homogeneous,
        $\Delta\Phi \approx 0$, and dark energy is negligibly small.
  \item In the late Universe---galaxies, clusters,
        superclusters---matter is strongly concentrated,
        $\Delta\Phi$ is large, and dark energy is significant.
\end{itemize}

\begin{hypothesis}[Coincidence problem]\label{hyp:coincidence}
In TGP the ``coincidental'' overlap between the epoch of
dark-energy domination and the epoch of structure formation
\textbf{is not a coincidence}---one causes the other.
Condensation of matter into structures generates additional
space, which we observe as dark energy.
\end{hypothesis}

% ----------------------------------------------------------
\subsection{Dark energy: natural $\Lambda$ or quintessence?}%
\label{ssec:de-dyn}

Numerical analysis (script \texttt{w\_de\_exact.py}) reveals
a central result.  For \textbf{natural} parameters
$\gamma \sim H_0^2/c_0^2$ (corresponding to
$\tau_0 \sim 1/H_0$, $\PhiZero \sim 25$), the field~$\varphi$
is \textbf{perfectly frozen}: Hubble friction $3H\dot\varphi$
dominates the potential force by $\sim 40$ orders of magnitude.
The outcome is $w_0 = -1.000$, $w_a = 0.000$---\textbf{indistinguishable
from a cosmological constant~$\Lambda$}.

TGP is therefore a \textbf{natural $\Lambda$~theory}: the
cosmological constant is not a parameter inserted by hand but
the residual potential $U(1) = \gamma/12$
(Proposition~\ref{prop:Lambda-eff}).

The effective cosmological constant is computable:
\begin{equation}\label{eq:Lambda-eff-def}
    \Lambda_{\mathrm{eff}}
    \;=\; \frac{8\pi G_0}{c_0^4}\;
    \bigl\langle U(\varphi_{\min}) \bigr\rangle_\Sigma,
\end{equation}
where $\varphi_{\min}(\mathbf{x})$ solves the field equation
subject to the zero-sum constraint, and
$\langle\cdot\rangle_\Sigma$ denotes the spatial average.

\medskip\noindent
\textbf{When is dark energy dynamical?}
The field~$\varphi$ evolves appreciably ($w_0 > -1$, $w_a < 0$)
only when $\gamma \gg H_0^2/c_0^2$, which requires either
unnaturally large~$\PhiZero$, higher-order substrate couplings,
or non-trivial initial conditions $\varphi_{\mathrm{ini}} \neq 1$.
In the dynamical regime TGP yields a \textbf{quintessence}-type
prediction:
\begin{itemize}[nosep]
  \item $w_0 \gtrsim -1$: today $\dot{\varphi}^2 > 0$ but small.
  \item $w_a < 0$: as structures stabilise,
        $w_{\mathrm{DE}} \to -1$ from above.
\end{itemize}

The equation of state of TGP dark energy takes the form
(Proposition~\ref{prop:wDE}):
\begin{equation}\label{eq:wDE-def}
    w_{\mathrm{DE}}
    \;=\; \frac{\frac{1}{2}\dot{\varphi}^2 - U(\varphi)}
               {\frac{1}{2}\dot{\varphi}^2 + U(\varphi)}.
\end{equation}
In the slow-relaxation regime
($\dot{\varphi}^2 \ll U(\varphi)$):
\begin{equation}\label{eq:wDE-slow}
    w_{\mathrm{DE}} \;\approx\; -1 + \delta w,
    \qquad
    \delta w \;=\; \frac{\dot{\varphi}^2}{U(\varphi)} \;>\; 0.
\end{equation}
For natural $\gamma \sim H_0^2/c_0^2$ one obtains
$\delta w < 10^{-40}$, operationally identical to $w = -1$.
For enhanced $\gamma$: $\delta w \sim 0.02$--$0.2$, yielding
testable quintessence with $(w_0 > -1,\; w_a < 0)$ in the CPL
parametrisation.

% ----------------------------------------------------------
\subsection{Relation to the zero-sum principle}\label{ssec:de-zero}

The zero-sum principle (Axiom~\ref{ax:zero}) requires
$\Ffield_{\mathrm{ext}} = 0$.  The additional space generated
by fluctuations \textbf{does not violate} this principle:
\begin{itemize}[nosep]
  \item An internal observer sees more space (dark energy).
  \item An external observer sees $\Ffield = 0$---the excess
        space is exactly compensated in the global balance.
\end{itemize}
Dark energy is therefore a \textbf{local self-organisation
effect of matter}, not a global constant.

\begin{remark}[Distinguishing TGP from $\Lambda$CDM]%
\label{rem:rozroznienie}
At natural parameters TGP is observationally indistinguishable
from $\Lambda$CDM via $w(z)$ measurements alone.
Discrimination requires \textbf{other observables}:
\begin{enumerate}[nosep]
  \item \textbf{GW polarisation}: a breathing mode alongside
        the tensorial $+/\times$ modes
        (Proposition~\ref{prop:GW-dispersion}).
  \item \textbf{Massive-GW dispersion relation}:
        a mass gap $f_{\mathrm{cut}} \sim H_0/(2\pi)$.
  \item \textbf{$\Lambda_{\mathrm{eff}}$ without fine-tuning}:
        $\Lambda \approx \gamma/12
        \sim \PhiZero\,H_0^2/(12\,c_0^2)$
        is \emph{calculable}, not inserted by hand.
  \item \textbf{Correlation of $\Lambda_{\mathrm{eff}}$
        with structure}: the back-reaction mechanism
        (Section~\ref{ssec:de-struktury}).
\end{enumerate}
\end{remark}
